\section{Đặt vấn đề}

Trong những năm gần đây, cùng với sự phát triển của thiết bị di động, mạng xã hội và các công cụ sáng tạo số, nhu cầu đăng tải, lưu trữ, chia sẻ và khai thác nội dung hình ảnh tăng lên rất nhanh. Ảnh không còn chỉ đóng vai trò minh họa cho văn bản, mà đã trở thành một loại nội dung trung tâm trong giao tiếp số, quảng bá cá nhân, xây dựng cộng đồng và phát triển kinh tế sáng tạo. Đối với các lĩnh vực như minh họa số, truyện tranh, thiết kế nhân vật, concept art, nhiếp ảnh, ảnh tạo bởi trí tuệ nhân tạo và các cộng đồng nghệ thuật trực tuyến, một nền tảng chia sẻ ảnh không chỉ là nơi lưu trữ tệp, mà còn là môi trường để người sáng tạo công bố tác phẩm, người xem khám phá nội dung, cộng đồng tương tác và các bên quản trị duy trì trật tự vận hành.

Sự phổ biến của các cộng đồng chia sẻ ảnh đặt ra một bài toán có quy mô lớn hơn so với các hệ thống quản lý ảnh truyền thống. Mỗi tác phẩm thường bao gồm nhiều thành phần: tệp ảnh gốc, phiên bản xem trước, ảnh thu nhỏ, mô tả, thẻ gắn kèm, thông tin tác giả, trạng thái xuất bản, quyền truy cập, bình luận, lượt thích, bộ sưu tập và báo cáo vi phạm. Khi số lượng người dùng và số lượng ảnh tăng, hệ thống phải xử lý đồng thời nhiều yêu cầu khác nhau: phục vụ hiển thị ảnh nhanh, kiểm soát nội dung không phù hợp, hỗ trợ tìm kiếm chính xác, bảo vệ quyền lợi của người sáng tạo và tạo ra cơ chế tương tác bền vững giữa tác giả với người theo dõi.

Một vấn đề cấp thiết trong các nền tảng chia sẻ ảnh là kiểm duyệt nội dung nhạy cảm hoặc độc hại. Với đặc thù cho phép người dùng tự do đăng tải, hệ thống có thể phải tiếp nhận các ảnh không phù hợp với chuẩn cộng đồng, vi phạm chính sách nội dung hoặc gây ảnh hưởng tiêu cực tới người xem. Nếu toàn bộ quá trình kiểm duyệt phụ thuộc vào con người, đội ngũ điều phối phải mở và đánh giá thủ công số lượng lớn ảnh, trong đó có cả nội dung nhạy cảm. Cách làm này tốn thời gian, khó mở rộng, tạo áp lực cho người kiểm duyệt và làm tăng độ trễ xuất bản tác phẩm. Ngược lại, nếu thiếu cơ chế kiểm soát, nền tảng có nguy cơ mất an toàn cộng đồng, giảm chất lượng nội dung và ảnh hưởng tới uy tín của người sáng tạo nghiêm túc.

Một vấn đề khác là phát hiện ảnh trùng lặp hoặc ảnh bị tái đăng trái phép. Trong môi trường nghệ thuật số, một tác phẩm có thể bị tải lại sau khi đã đổi kích thước, nén lại, cắt nhẹ hoặc thay đổi định dạng. Việc chỉ so sánh tên tệp, kích thước tệp hoặc mã băm nhị phân không đủ để phát hiện các trường hợp gần giống nhau về mặt thị giác. Nếu bài toán này không được xử lý, nền tảng dễ bị lặp nội dung, kết quả tìm kiếm bị nhiễu, quyền lợi của tác giả bị ảnh hưởng và công việc kiểm duyệt bản quyền trở nên khó khăn hơn. Đây là vấn đề có ý nghĩa thực tế lớn đối với các cộng đồng sáng tạo, nơi tính nguyên bản của tác phẩm là yếu tố quan trọng.

Bên cạnh đó, tìm kiếm ảnh dựa trên từ khóa còn nhiều hạn chế. Trên nhiều nền tảng, người dùng chỉ có thể tìm kiếm dựa vào tiêu đề, mô tả hoặc tag do tác giả nhập. Cách tìm kiếm này phụ thuộc mạnh vào chất lượng mô tả văn bản và sự trùng khớp giữa từ khóa của người tìm với từ khóa của người đăng. Trong khi đó, nội dung thị giác của ảnh thường chứa những đặc trưng khó mô tả bằng vài từ ngắn, chẳng hạn bố cục, phong cách nét vẽ, màu sắc chủ đạo, cảm xúc nhân vật hoặc mức độ tương đồng giữa một bản phác thảo và một tác phẩm hoàn thiện. Khi người dùng không biết chính xác tag cần nhập, hoặc khi tác giả gắn tag chưa đầy đủ, khả năng khám phá nội dung phù hợp bị giảm đáng kể.

Cuối cùng, các nền tảng chia sẻ ảnh hiện đại cần quan tâm tới nhu cầu tạo thu nhập của người sáng tạo. Nhiều nghệ sĩ không chỉ muốn đăng tải tác phẩm công khai, mà còn muốn cung cấp nội dung độc quyền, ảnh độ phân giải cao, bản phác thảo, tài liệu hậu trường hoặc quyền truy cập sớm cho nhóm người hâm mộ trả phí. Nếu nền tảng chia sẻ ảnh không hỗ trợ cơ chế thành viên, người sáng tạo phải sử dụng thêm các dịch vụ bên ngoài để quản lý đăng ký và phân phối nội dung riêng. Điều này làm phân mảnh trải nghiệm người dùng, giảm khả năng giữ chân cộng đồng trên cùng một nền tảng và khiến việc kiểm soát quyền truy cập nội dung trở nên phức tạp.

Từ các vấn đề trên, có thể thấy bài toán xây dựng nền tảng chia sẻ ảnh trực tuyến có ý nghĩa thực tiễn rõ ràng. Nếu được giải quyết tốt, hệ thống không chỉ hỗ trợ người dùng phổ thông trong việc tìm kiếm và tương tác với nội dung ảnh, mà còn hỗ trợ người sáng tạo bảo vệ tác phẩm, quản lý cộng đồng người theo dõi và phát triển nguồn thu trực tiếp. Các kết quả của bài toán cũng có thể được mở rộng cho những miền ứng dụng khác như thư viện ảnh doanh nghiệp, nền tảng học tập mỹ thuật, kho dữ liệu thiết kế, hệ thống quản lý nội dung truyền thông hoặc các cộng đồng sáng tạo có dữ liệu đa phương tiện.

\section{Mục tiêu và phạm vi đề tài}

Hiện nay đã tồn tại nhiều sản phẩm phục vụ một phần nhu cầu của bài toán chia sẻ ảnh. Các mạng xã hội hình ảnh phổ thông có ưu điểm về khả năng lan truyền nội dung và tương tác nhanh, nhưng thường không tập trung vào quy trình xử lý ảnh chuyên sâu, phát hiện trùng lặp hoặc kiểm duyệt theo đặc thù của cộng đồng nghệ thuật. Các cộng đồng nghệ thuật số như các nền tảng hồ sơ nghệ sĩ hỗ trợ đăng tác phẩm, tag và bình luận, song tìm kiếm vẫn chủ yếu dựa vào từ khóa và mức độ đầy đủ của mô tả do người đăng nhập. Các nền tảng đăng ký hội viên hỗ trợ người sáng tạo thu phí từ người hâm mộ, nhưng thường tách rời khỏi nơi người dùng khám phá và tương tác với tác phẩm. Một số công cụ tìm kiếm ảnh bằng trí tuệ nhân tạo có khả năng hiểu nội dung thị giác tốt hơn, nhưng lại không bao phủ đầy đủ các chức năng xã hội, kiểm duyệt, phân quyền và thanh toán.

Như vậy, hạn chế chung của các hướng tiếp cận hiện tại là sự phân tán chức năng. Người sáng tạo có thể cần một nơi để đăng tác phẩm, một nơi khác để nhận thành viên trả phí, một công cụ riêng để quản lý ảnh và một quy trình riêng để kiểm duyệt hoặc xử lý tranh chấp. Đối với người xem, trải nghiệm khám phá cũng bị hạn chế khi tìm kiếm chỉ dựa trên từ khóa và không tận dụng được đặc trưng thị giác của ảnh. Đối với quản trị viên, việc kiểm duyệt thủ công trên khối lượng ảnh lớn có thể gây quá tải. Từ đó, đồ án hướng tới xây dựng một nền tảng web chia sẻ ảnh trực tuyến tích hợp nhiều nhóm chức năng cần thiết trong cùng một hệ thống.

Mục tiêu tổng quát của đề tài là phát triển một nền tảng chia sẻ ảnh trực tuyến phục vụ cộng đồng nghệ thuật số, trong đó người dùng có thể khám phá, tìm kiếm, xem và tương tác với tác phẩm; người sáng tạo có thể đăng tải, quản lý tác phẩm và thiết lập gói thành viên; quản trị viên có thể kiểm duyệt nội dung, xử lý báo cáo và theo dõi hoạt động của hệ thống. Hệ thống được xây dựng dưới dạng ứng dụng web, với phần giao diện dành cho khách, người dùng đã đăng nhập, người sáng tạo và nhóm quản trị.

Về phạm vi chức năng, hệ thống tập trung vào các nhóm chức năng chính. Nhóm thứ nhất là quản lý tài khoản và phân quyền, bao gồm xác thực người dùng, đồng bộ hồ sơ, phân biệt người dùng thường, người sáng tạo và quản trị viên. Nhóm thứ hai là quản lý tác phẩm, bao gồm đăng tải nhiều ảnh, lưu thông tin mô tả, quản lý trạng thái xử lý, hỗ trợ chèn dấu bản quyền và kiểm soát quyền truy cập nội dung. Nhóm thứ ba là tìm kiếm và khám phá, bao gồm tìm kiếm từ khóa, lọc theo thuộc tính, xem nội dung nổi bật, tìm kiếm ngữ nghĩa và gợi ý tác phẩm liên quan. Nhóm thứ tư là tương tác cộng đồng, bao gồm thích, bình luận, lưu vào bộ sưu tập, theo dõi nghệ sĩ và báo cáo vi phạm. Nhóm thứ năm là cơ chế thành viên và thanh toán, bao gồm gói thành viên của người sáng tạo, đăng ký, hủy đăng ký, lịch sử thanh toán, doanh thu và thanh toán. Nhóm cuối cùng là quản trị và kiểm duyệt, bao gồm xử lý báo cáo, duyệt hoặc từ chối tác phẩm, quản lý người dùng, cảnh báo, khóa tài khoản và xem nhật ký thao tác.

Về phạm vi kỹ thuật, đề tài không đặt mục tiêu huấn luyện lại mô hình trí tuệ nhân tạo từ đầu hoặc xây dựng hạ tầng phân tán quy mô rất lớn. Các mô hình và thư viện có sẵn được sử dụng như thành phần hỗ trợ trong hệ thống ứng dụng. Đồ án cũng không tập trung vào các vấn đề ngoài phạm vi như thương mại hóa sản phẩm, tối ưu hạ tầng đa vùng, xây dựng ứng dụng di động native hoặc giải quyết toàn bộ quy trình pháp lý về bản quyền. Trọng tâm của đề tài là thiết kế và triển khai một nền tảng web có kiến trúc rõ ràng, có khả năng xử lý ảnh bất đồng bộ, hỗ trợ tìm kiếm đa dạng, có kiểm duyệt và có mô hình thành viên cho người sáng tạo.

\section{Định hướng giải pháp}

Từ mục tiêu và phạm vi trên, đồ án định hướng xây dựng hệ thống theo kiến trúc modular monolith. Theo định hướng này, backend được triển khai như một ứng dụng thống nhất nhưng được chia thành các module nghiệp vụ rõ ràng như xác thực, người dùng, tác phẩm, tìm kiếm, hàng đợi xử lý, thanh toán, báo cáo, quản trị, khuyến nghị và theo dõi. Cách tiếp cận này phù hợp với quy mô của đồ án vì giữ được sự đơn giản trong triển khai, đồng thời vẫn tạo ranh giới mã nguồn đủ rõ để mở rộng hoặc tách dịch vụ trong tương lai.

Đối với xử lý ảnh và các tác vụ tốn tài nguyên, hệ thống sử dụng mô hình xử lý nền thông qua hàng đợi. Khi người sáng tạo đăng tải tác phẩm, các bước nặng như tạo ảnh thu nhỏ, tạo ảnh làm mờ, kiểm tra trùng lặp, kiểm duyệt nội dung và sinh đặc trưng ảnh không được thực hiện trực tiếp trong luồng yêu cầu chính. Định hướng này giúp giảm thời gian chờ của người dùng, hạn chế lỗi quá thời gian chờ và cho phép hệ thống theo dõi trạng thái xử lý của từng tác phẩm.

Đối với tìm kiếm, hệ thống kết hợp tìm kiếm toàn văn và tìm kiếm ngữ nghĩa. Tìm kiếm toàn văn phục vụ các truy vấn dựa trên tiêu đề, mô tả, tag và tên tác giả; tìm kiếm ngữ nghĩa phục vụ các truy vấn tự nhiên hoặc truy vấn dựa trên đặc trưng thị giác của ảnh. Việc kết hợp hai hướng này giúp hệ thống vừa đáp ứng nhu cầu tìm kiếm truyền thống, vừa cải thiện khả năng khám phá nội dung trong trường hợp từ khóa không đủ diễn đạt nội dung ảnh.

Đối với kiểm duyệt và bảo vệ nội dung, hệ thống định hướng sử dụng các kỹ thuật tự động ở bước đầu và vẫn giữ vai trò xử lý cuối cùng cho quản trị viên. Nội dung nhạy cảm được phát hiện bằng mô hình phân loại ảnh; ảnh trùng lặp hoặc gần trùng lặp được phát hiện bằng dấu vân tay thị giác; các trường hợp bị báo cáo hoặc cần xem xét lại được đưa vào quy trình kiểm duyệt. Cách tiếp cận này giúp giảm tải cho người điều phối nhưng không loại bỏ hoàn toàn yếu tố con người trong các quyết định nhạy cảm.

Đối với mô hình kinh tế sáng tạo, hệ thống tích hợp cơ chế gói thành viên của người sáng tạo và thanh toán trực tuyến. Người sáng tạo có thể tạo các gói thành viên với mức giá và quyền lợi khác nhau; người dùng có thể đăng ký để truy cập nội dung giới hạn; hệ thống đồng bộ trạng thái thanh toán để kiểm soát quyền truy cập. Đây là thành phần quan trọng để nền tảng không chỉ là nơi chia sẻ ảnh mà còn hỗ trợ trực tiếp hoạt động sáng tạo chuyên nghiệp.

Đóng góp chính của đồ án là xây dựng một nền tảng chia sẻ ảnh trực tuyến tích hợp các nhóm chức năng thường bị tách rời: đăng tải và xử lý ảnh, kiểm duyệt bán tự động, phát hiện trùng lặp, tìm kiếm lai, tương tác cộng đồng, quản lý thành viên và thanh toán. Kết quả đạt được là một hệ thống web có cấu trúc module, có luồng xử lý nền, có mô hình dữ liệu phục vụ ảnh và tác phẩm, đồng thời có nền tảng để mở rộng sang các chức năng nâng cao như khuyến nghị cá nhân hóa, kiểm duyệt nâng cao và phân tích cộng đồng.

\section{Bố cục đồ án}

Phần còn lại của báo cáo đồ án tốt nghiệp được tổ chức thành năm chương tiếp theo. Chương 2 trình bày khảo sát và phân tích yêu cầu của hệ thống. Nội dung chương này bao gồm khảo sát hiện trạng các nhóm nền tảng chia sẻ ảnh và các hệ thống liên quan, xác định các tác nhân tham gia, mô tả các chức năng mức cao, xây dựng biểu đồ use case, đặc tả một số use case quan trọng và tổng hợp các yêu cầu phi chức năng. Đây là cơ sở để xác định phạm vi thiết kế và triển khai ở các chương sau.

Chương 3 trình bày nền tảng lý thuyết và công nghệ sử dụng trong đồ án. Chương này giới thiệu các công nghệ nền tảng cho backend, frontend, cơ sở dữ liệu, xử lý nền, tìm kiếm, xác thực và thanh toán. Bên cạnh đó, chương cũng trình bày các cơ sở lý thuyết liên quan tới mô hình học biểu diễn ảnh-văn bản, phân loại ảnh nhạy cảm, băm nhận thức và tìm kiếm trong không gian vector. Các nội dung này giúp giải thích vì sao hệ thống lựa chọn những kỹ thuật và công nghệ tương ứng.

Chương 4 trình bày quá trình phân tích thiết kế, triển khai và đánh giá hệ thống. Nội dung chương bao gồm thiết kế kiến trúc tổng thể, thiết kế cơ sở dữ liệu, thiết kế các module backend, thiết kế giao diện chính, mô tả luồng xử lý ảnh, cấu hình các dịch vụ trung gian và phương pháp kiểm thử. Chương này làm rõ cách các yêu cầu đã nêu ở Chương 2 được chuyển hóa thành thành phần kỹ thuật cụ thể.

Chương 5 trình bày các giải pháp và đóng góp nổi bật của đồ án. Chương này tập trung vào những điểm có giá trị kỹ thuật và nghiệp vụ của hệ thống, bao gồm tách xử lý ảnh khỏi luồng yêu cầu chính, kiểm duyệt bán tự động, phát hiện ảnh trùng lặp, tìm kiếm ngữ nghĩa, đồng bộ chỉ mục tìm kiếm và mô hình thành viên cho người sáng tạo. Các nội dung trong chương này làm rõ mức độ đóng góp của đồ án so với một hệ thống chia sẻ ảnh thông thường.

Chương 6 đưa ra kết luận và hướng phát triển. Chương này tổng kết các kết quả đã đạt được, đánh giá mức độ đáp ứng mục tiêu ban đầu, nêu các hạn chế còn tồn tại và đề xuất các hướng mở rộng trong tương lai. Những hướng phát triển có thể bao gồm nâng cao chất lượng kiểm duyệt, tối ưu tìm kiếm vector, cải tiến hệ thống khuyến nghị, bổ sung cơ chế xử lý tranh chấp bản quyền và mở rộng hạ tầng triển khai cho quy mô người dùng lớn hơn.
